\section{Executive conclusion}
\label{sec:exec-conclusion}

BayesBreak addresses a well-defined statistical problem: Bayesian multiple-changepoint segmentation with observation-family generality and hierarchical structure for irregular designs, multiple sequences, and observed or latent groups. Its principal mathematical mechanism is standard in form but broad in use: integrate segment parameters, combine segment marginal likelihoods with a factorized partition prior, and use dynamic programming for posterior marginalization or joint MAP optimization. The contribution lies in the generalized hierarchical formulation, the exact results under stated conditions, the explicit distinction between exact and approximate segment calculations, and the implementation across several observation families and structured settings.

The technical book and journal paper establish the current method without changing its title or research direction. The archived numerical values are real executed results and remain read-only. The next work is not a conceptual rebranding exercise. It is a defined program of statistical implementation, numerical verification, corrected computation, repeated validation, and external release review.

The immediate decision is whether to fund only the minimum corrective cycle through Stage S5 or the broader validation and release program through Stage S7. In either case, claims must remain matched to proofs, probability models, optimization objectives, and executed experiments. No invalid comparison, unsupported approximation rate, or unresolved test state should be converted into a stronger conclusion through wording alone.
